% ============================================
% Johns Hopkins Letter Template
% Assistant Professor Cover Letter – UMass Boston
% ============================================

\documentclass[11pt,letterpaper]{letter}

\usepackage{graphicx}
\usepackage{eso-pic}
\usepackage{geometry}
\usepackage{microtype}
\usepackage[hidelinks]{hyperref}

% ----- Page Setup -----
\geometry{left=25mm,right=25mm,top=25mm,bottom=25mm}

% ----- Logo Placement (first page only) -----
\newcommand{\PutJHULogoFG}[1]{%
  \AddToShipoutPictureFG*{%
    \AtPageUpperLeft{%
      \hspace{10mm}%
      \raisebox{-38mm}{%
        \includegraphics[width=6cm]{#1}%
      }%
    }%
  }%
}

% ----- Sender -----
\signature{Decory Edwards}
\address{Department of Economics \\ Johns Hopkins University \\ Baltimore, MD 21218 \\ \texttt{dedwar65@jh.edu}}

\begin{document}
\PutJHULogoFG{Images/jhulogo.png}

\begin{letter}{Search Committee \\
Department of Economics \\
University of Massachusetts Boston \\
100 Morrissey Boulevard \\
Boston, MA 02125 \\ USA}

\opening{Dear Members of the Search Committee,}

I am writing to express my interest in the tenure-track Assistant Professor position in the Department of Economics at the University of Massachusetts Boston, beginning September 2026. I am a Ph.D. candidate in Economics at Johns Hopkins University, advised by Professors Christopher D. Carroll, Nicholas W. Papageorge, and Stelios Fourakis, and I expect to receive my degree in May 2026. My research fields are macroeconomics, wealth and income dynamics, and computational economics.

My research agenda focuses on understanding the determinants of wealth inequality and the mechanisms through which heterogeneity in returns to wealth shapes aggregate outcomes. In my job-market paper, I estimate the degree of heterogeneity in individual portfolio returns using micro-data from the Survey of Consumer Finances and simulate a structural model in which households face idiosyncratic return risk. I show that allowing for persistent heterogeneity in returns substantially improves the model’s ability to match the empirical distribution of wealth, offering a tractable way to connect micro-level return dispersion to macroeconomic inequality.

In a second project, I use the Health and Retirement Study to examine the relationship between interpersonal trust and portfolio performance. Building on the literature that finds a hump-shaped relationship between trust and income, I show that a similar relationship holds for individual returns to wealth measured in the HRS data, highlighting a behavioral channel through which social attitudes shape saving and investment outcomes. This work also provides new estimates of the persistent component of returns to net worth, which motivate the calibration of heterogeneity in my structural model.

UMass Boston’s Department of Economics emphasizes rigorous empirical research, policy relevance, and a strong commitment to addressing inequality and economic justice—an agenda that aligns closely with my work. The university’s role as a major public research institution serving a diverse urban population is especially appealing to me, and I see strong synergies between my research on wealth dynamics, heterogeneous returns, and macroeconomic inequality and the department’s strengths in applied policy analysis and econometrics. I am excited about contributing to an environment where research informs public policy and where diverse student backgrounds enrich the study of economics.

\textbf{Teaching.} My teaching experience centers on undergraduate macroeconomics and an upper-level macro policy seminar, where I have led weekly discussion sections and held regular office hours for classes ranging from 16 to 40 students. My approach is student-centered: I aim to help students “convince themselves” that they have moved from confusion to understanding by holding structured office-hour coaching that builds trust and confidence. At UMass Boston, I am prepared to teach \textit{Principles of Macroeconomics}, \textit{Intermediate Macroeconomic Theory}, \textit{Money and Banking}, and \textit{Quantitative Methods for Economics}. I would also welcome developing electives such as \textit{Economics of Inequality and Tax Policy} or \textit{Computational Macroeconomics and Policy}, emphasizing reproducible workflows, transparent coding practices, and evidence-based decision-making.

I am committed to fostering an inclusive learning environment consistent with UMass Boston’s urban mission and its diverse student body. In my teaching and mentorship, I incorporate accessible quantitative tools, a variety of assessment formats, and research opportunities that support students with differing backgrounds and preparation. I would look forward to contributing to the department’s teaching, research, and service responsibilities while supporting a student community that brings a wide range of experiences to the study of economics.

I would be honored to join the University of Massachusetts Boston’s Department of Economics and contribute to its teaching and research missions. Thank you for your consideration; I welcome the opportunity to discuss my fit with the department at your convenience.

\closing{Sincerely,}

\end{letter}
\end{document}
